% !TeX root = ../../../main.tex
\section{Discussion}
\label{sec:discussion}

\subsection{Insights on Times-Intersect and its Measurement}

Fig.~\ref{fig:tintDistrib} shows three important trends. First, across all beam sequences, $t_{\textnormal{int}}$ decreased nonlinearly with increasing shear modulus. For a given change in modulus, differences in $t_{\textnormal{int}}$ were greater for a pair of softer materials versus those of stiffer materials. This trend is important because, since DoPIo estimates modulus from $t_{\textnormal{int}}$, it implies that distinguishing shear elastic moduli in stiffer media is inherently more challenging than among lower ones.

Second, when separate tracking is performed, the dynamic range of $t_{\textnormal{int}}$, as well as differences in their magnitudes between moduli, were greater when the ARF push beam was wider and the difference in tracking beam widths was larger. In particular, differences in $t_{\textnormal{int}}$ between moduli were largest using the beam sequence that used an F/3.0 ARF excitation and F/1.5 and F/5.0 tracking. Relative to the other examined beam sequences, this beam sequence also achieved the largest effect sizes, as shown in Figure~\ref{fig:tintStat}. This trend aligns with the conceptual basis of DoPIo's approach, where shear elastic modulus is estimated by leveraging scatterer shearing. Faster shearing rates are more effectively distinguished when differences in $t_{\textnormal{int}}$ are greater. This can occur when (1) the difference in lateral widths of the two tracking beams is larger, as well as (2) when scatterers displace in distinctly different manners under one beam's PSF versus the other's. This phenomenon was observed when a F/1.5 and F/5.0 track beam combination was used, as those beams had a greater difference in lateral widths compared to the F/1.5 and F/3.0 pair. Paired with a wider F/3.0 ARF push, the F/1.5 track beam primarily captured scatterers near the center of the ARF focus that displaced relatively uniformly at the time of excitation. Meanwhile, the F/5.0 beam applies a heavier weighting to scatterers located beyond the main lobe of the ARF push beam. Those outer scatterers displace later, at a time influenced by the rate of scatterer shearing, leading to a more distinct displacement profile that enhances the measurable difference in $t_{\textnormal{int}}$.

Third, for that sequence involving an F/3.0 push and F/1.5 and F/5.0 tracking, separate and simultaneous tracking yielded comparable $t_{\textnormal{int}}$ distributions (Fig.~\ref{fig:tintDistrib}). Simultaneous tracking offers advantages in DoPIo by doubling the frame rate, halving the number of required ARF excitations, and reducing susceptibility to motion artifacts. However, it requires transmitting a single pulse with a width matching the wider of the two tracking beams, followed by beamforming the resulting channel data to achieve two distinct receive focal configurations. A potential drawback of this approach is the reduced difference in beam widths, which can lead to smaller $t_{\textnormal{int}}$ differences and, consequently, diminished accuracy of modulus estimation - particularly for stiffer materials. This limitation was most evident for the beam sequence using an F/3.0 ARF excitation with F/1.5 and F/3.0 tracking. Since both tracking beams spanned scatterers that were displaced at the time of the ARF excitation and the difference in beam widths was smaller due to the use of simultaneous tracking, stiffness-based trends in $t_{\textnormal{int}}$ were less clear than in the separate tracking case. However, for the beam sequencing using an F/3.0 ARF excitation and F/1.5 and F/5.0 tracking, $t_{\textnormal{int}}$ distributions across elasticities were similar between separate and simultaneous tracking despite the smaller difference in tracking beam widths for the latter case. 

\subsection{Relating Times-Intersect to Shear Elastic Modulus}

The data included in Table~\ref{tbl:modelr2} demonstrate that the arbitrary power law model (Eq.~\ref{eq:modelarb}) consistently yielded higher coefficients of determination compared to the inverse-square model. However, the derived powers varied from 2.1 at the 25-mm focal depth to 1.8 within 10 mm above or below the focus, closely adhering to the inverse-square relationship expected between shear wave propagation time and shear elastic modulus~\cite{sarvazyan_Shear_1998,palmeri_Ultrasonic_2006}. The deviation from an inverse-square relationship could attribute for the effects of the ARF spatial gradient and duration as well as compression waves, mode conversion, and other factors that impact scatterer shear rate before a planar shear wave is fully formed~\cite{landau_Theory_1986,sarvazyan_Shear_1998,ostrovsky_NonDissipative_2006}.

When using the arbitrary power law model, Fig.~\ref{fig:Gdistrib} and Table~\ref{tbl:medianG} show that DoPIo's shear elastic modulus estimates were generally accurate across all beam sequences, with median errors smaller than the simulated modulus step size of 4 kPa. Notably, the sequence combining an F/3.0 ARF excitation with F/1.5 and F/5.0 simultaneous tracking produced the most accurate estimates, consistent with its ability to generate the most distinct $t_{\textnormal{int}}$ distributions. However, modulus estimates from all evaluated beam sequences exhibited lower precision in stiffer materials, as indicated by larger interquartile ranges (IQRs) in the modulus estimates in Fig.~\ref{fig:Gdistrib}. However, this reduced precision is expected since $t_{\textnormal{int}}$ distributions among stiffer materials tend to overlap more than those among softer materials, making them harder to distinguish. While this limitation may be partially addressed by further optimizing $t_{\textnormal{int}}$ detection through improved DoPIo beamforming or increased PRF, an alternative modeling approach may ultimately be required to enhance precision in stiffer materials. Such an approach could include machine learning techniques that directly leverage paired displacement profiles as inputs, rather than relying solely on $t_{\textnormal{int}}$ as an intermediate metric. Initial results using a support vector machine-based regression model have shown promise in improving modulus estimation precision~\cite{rahmouni_Machine_2020}, though further investigation using more advanced techniques is warranted.

\subsection{\emph{In Vivo} and \emph{Ex Vivo} Acquisitions}

Fig.~\ref{fig:cirsImage} shows that elastic inclusions within the calibrated phantom were successfully delineated in the DoPIo images. Qualitatively compared to the background, the 2.2 and 5.1 kPa inclusions exhibited longer $t_{\textnormal{int}}$ and lower DoPIo-estimated elastic moduli while the 16.3 kPa inclusion showed shorter times-intersect and higher estimated moduli. Quantitatively, DoPIo tended to underestimate the modulus of most features, with median and median absolute deviation errors being larger for stiffer materials (e.g. -3.3 ± 0.5 kPa in the 16.3 kPa inclusion) versus softer ones (-0.6 ± 0.2 kPa for the 5.1 kPa inclusion). This discrepancy may stem from increased overlap in $t_{\textnormal{int}}$ distributions in stiffer materials, leading to less accurate and less precise modulus estimates. Another contributing factor could be limitations in the empirical model used to relate $t_{\textnormal{int}}$ to modulus. Since this model was derived \emph{in silico}, it may not be optimally calibrated for \emph{in vitro} applications. This mismatch is potentially due to inconsistencies between simulated and acquired displacement data or greater measurement error in $t_{\textnormal{int}}$ during experimental acquisitions. The latter issue may be particularly pronounced in stiffer materials, which undergo smaller displacements in response to ARF excitation, resulting in displacement magnitudes that approach the inherent variance of the normalized cross-correlation displacement estimator~\cite{mcaleavey_Estimates_2003,dhanaliwala_Assessing_2012}.

Fig.~\ref{fig:cirsBoxplot} investigates the influence of ARF-induced displacement magnitude on DoPIo modulus estimates. Panels (a)-(f) show that the estimated modulus in the 5.1 kPa inclusion remained stable as the ARF voltage increased from 35 V to 60 V. In contrast, the estimated modulus in the 8.7 kPa background rose between 35 V and 50 V, then plateaued from 50 V to 60 V. These trends are echoed in the box plots in panel~(g). DoPIo modulus estimates for the 2.2 kPa inclusion were consistent across all ARF voltage levels. However, for the 16.3 and 29.0 kPa inclusions, modulus estimates increased with ARF voltage from 35 V to 55 V, then stabilized at 55 V and 60 V. These findings suggest that a minimum threshold of ARF-induced displacement is necessary for reliable times-intersect measurements. Furthermore, for a given focal depth and DoPIo beam sequence, stiffer materials require higher ARF push amplitudes to reach this threshold. Importantly, once sufficient displacement is achieved, DoPIo modulus estimates appear to be independent of ARF amplitude, reinforcing DoPIo's potential for quantitative, on-axis, ARF-based elasticity imaging.

Panel (g) also reveals a noteworthy discrepancy: even at the highest ARF push voltage, DoPIo underestimated the modulus of the 29.0 kPa inclusion by approximately 50\%, exceeding the error observed \emph{in silico} in Fig.~\ref{fig:Gdistrib}. This unexpected deviation may result from miscalibration of the empirical $t_{\textnormal{int}}$-to-modulus model for \emph{in vitro} conditions as previously discussed. Alternatively, it may reflect limitations in accurately and precisely measuring low times-intersect in stiff materials during experimental acquisition. To address this, strategies to increase times-intersect, such as enhancing the difference in lateral widths of the paired tracking beams or increasing the PRF, could be explored. However, these adjustments would come with trade-offs: increasing the difference in lateral beam widths would reduce the lateral field of view for 2D DoPIo imaging and degrade radiofrequency signal-to-noise ratio, while the use of higher PRFs would limit imaging depth.

Fig.~\ref{fig:liver} presents DoPIo imaging results from excised liver tissue. As shown in panels (b) and (c), the stiffer gelatin inclusion exhibited lower times-intersect and higher estimated modulus compared to the softer liver background. Panel~(d) displays boxplots of DoPIo modulus estimates across all examined ARF voltage levels for both materials. In general, modulus estimates remained consistent across voltages for both gelatin and liver features, indicating that sufficient ARF-induced displacements were achieved to capture reliable times-intersect measurements. However, modulus estimates were consistently lower than those derived from clinical scanner-based validation standards. This underestimation conflicts with DoPIo's performance in the calibrated phantom, where modulus estimates for media with sub-10 kPa shear moduli were accurate. One possible explanation is the viscoelastic nature of liver tissue. Viscosity may dampen scatterer shearing, leading to longer times-intersect and, consequently, lower estimated moduli. Future work will focus on developing methods to account for, and possibly directly interrogate, viscoelastic effects in DoPIo imaging. To achieve this goal, modified acquisition parameters may be employed to enhance DoPIo's sensitivity to viscous effects, and new data-driven approaches may be explored to exploit viscoelastic behavior that is encoded in displacement profiles.
